\documentclass{article}

% Numbered, compressed citations: \cite{a} -> [1], \cite{a,b,c} -> [1-3].
\PassOptionsToPackage{numbers, compress}{natbib}

% The authors should use one of these tracks.
% Before accepting by the NeurIPS conference, select one of the options below.
% 0. "default" for submission
\usepackage{neurips_2026}
% the "default" option is equal to the "main" option, which is used for the Main Track with double-blind reviewing.
% 1. "main" option is used for the Main Track
%  \usepackage[main]{neurips_2026}
% 2. "position" option is used for the Position Paper Track
%  \usepackage[position]{neurips_2026}
% 3. "eandd" option is used for the Evaluations & Datasets Track
 % \usepackage[eandd]{neurips_2026}
 % if you need to opt-in for a single-blind submission in the E&D track:
 %\usepackage[eandd, nonanonymous]{neurips_2026}
% 4. "creativeai" option is used for the Creative AI Track
%  \usepackage[creativeai]{neurips_2026}
% 5. "sglblindworkshop" option is used for the Workshop with single-blind reviewing
 % \usepackage[sglblindworkshop]{neurips_2026}
% 6. "dblblindworkshop" option is used for the Workshop with double-blind reviewing
%  \usepackage[dblblindworkshop]{neurips_2026}

% After being accepted, the authors should add "final" behind the track to compile a camera-ready version.
% 1. Main Track
 % \usepackage[main, final]{neurips_2026}
% 2. Position Paper Track
%  \usepackage[position, final]{neurips_2026}
% 3. Evaluations & Datasets Track
 % \usepackage[eandd, final]{neurips_2026}
% 4. Creative AI Track
%  \usepackage[creativeai, final]{neurips_2026}
% 5. Workshop with single-blind reviewing
%  \usepackage[sglblindworkshop, final]{neurips_2026}
% 6. Workshop with double-blind reviewing
%  \usepackage[dblblindworkshop, final]{neurips_2026}
% Note. For the workshop paper template, both \title{} and \workshoptitle{} are required, with the former indicating the paper title shown in the title and the latter indicating the workshop title displayed in the footnote.
% For workshops (5., 6.), the authors should add the name of the workshop, "\workshoptitle" command is used to set the workshop title.
% \workshoptitle{WORKSHOP TITLE}

% "preprint" option is used for arXiv or other preprint submissions
 % \usepackage[preprint]{neurips_2026}

% to avoid loading the natbib package, add option nonatbib:
%    \usepackage[nonatbib]{neurips_2026}

\usepackage[utf8]{inputenc} % allow utf-8 input
\usepackage[T1]{fontenc}    % use 8-bit T1 fonts
\usepackage{hyperref}       % hyperlinks
\usepackage{url}            % simple URL typesetting
\usepackage{booktabs}       % professional-quality tables
\usepackage{amsfonts}       % blackboard math symbols
\usepackage{nicefrac}       % compact symbols for 1/2, etc.
\usepackage{microtype}      % microtypography
\usepackage{xcolor}         % colors
\usepackage{tabularx}
\usepackage{amsmath}
  \usepackage{mathtools}
 \usepackage{tikz}
  \usetikzlibrary{positioning,arrows.meta,fit,backgrounds}
\usepackage{longtable}

% Note. For the workshop paper template, both \title{} and \workshoptitle{} are required, with the former indicating the paper title shown in the title and the latter indicating the workshop title displayed in the footnote. 
\title{A coloring-free approach to materialising Jacobians and Hessians in JAX by leveraging sparsity}


% The \author macro works with any number of authors. There are two commands
% used to separate the names and addresses of multiple authors: \And and \AND.
%
% Using \And between authors leaves it to LaTeX to determine where to break the
% lines. Using \AND forces a line break at that point. So, if LaTeX puts 3 of 4
% authors names on the first line, and the last on the second line, try using
% \AND instead of \And before the third author name.


\author{%
  Jonathan Brodrick\\
  Pasteur Labs\\
  Newlab, 19 Morris Street\\
  Brooklyn, NY 11205 \\
  \texttt{jonathan.brodrick@simulation.science} \\
  % examples of more authors
  % \And
  % Coauthor \\
  % Affiliation \\
  % Address \\
  % \texttt{email} \\
  % \AND
  % Coauthor \\
  % Affiliation \\
  % Address \\
  % \texttt{email} \\
  % \And
  % Coauthor \\
  % Affiliation \\
  % Address \\
  % \texttt{email} \\
  % \And
  % Coauthor \\
  % Affiliation \\
  % Address \\
  % \texttt{email} \\
}


\begin{document}


\maketitle


\begin{abstract}
Automatic differentiation (AD) has proven indispensable in powering the ML revolution by enabling accelerated
gradient-based learning without hard-to-maintain hand-coded analytic derivatives.
However, modern AD libraries such as JAX still leave significant room for improvement when it comes to computation of 
full Jacobians and Hessians which are central to root finding (e.g. in implicit ODE solvers or deep equilibrium models)
or second-order constrained optimisation (relevant to robotic and optimal control).
In this work, we show that leveraging JAX's functional transformations and introducing new sparse formats,
including batched Ellpack, can reduce redundant computation and exploit internal structure achieving orders of 
magnitude speedups in structured nonlinear problems (including the CUTEst suite).
This functionality is made available in an open source library with a drop-in API mirroring JAX's
(including hessian, jacfwd and jacrev with a choice of output formmats).
Our methodology is competitive with colouring methods and outperforms it on some problems
without requiring the user to manually provide the sparsity pattern 
and avoiding the sometimes expensive coloring phase.
Furthermore, when the sparsity pattern is statically known it can be obtained for free at trace time.
We envisage the key use case to be coupling the accelerated Jacobian/Hessian computation to a sparse solver such as 
CudSS.
\end{abstract}doe


\section{Introduction}

The JAX library provides a powerful framework for implementing large ML models/simulators 
and training/optimizing them using gradient-based methods (e.g. Adam, L-BFGS) on modern accelerators.
For these cases, the computational efficiency of the gradient as obtained by JAX's automatic differentiation will usually be competitive with
any handcoded version, greatly reducing development time and improving maintainability.
However, the same cannot be said for the case of root finding (e.g. in implicit ODE solvers or deep equilibrium models) 
 and second-order optimisation using direct solvers (relevant to robotics and optimal control)
 where the entire Jacobian and Hessian must be obtained.
This is because JAX computes these simply by applying the jvp/hvp to a batch of basis vectors rather than 
leveraging structural sparsity;
the resulting $\mathcal{O}(n^2)$ cost of computing the Jacobian/Hessian is most acute when 
the matrices possess structure (e.g. sparsity) that enables 
Traditional ways . 
function transformation


\section{State of the art}

section{Methodology}

\subsection{Sparse representations}

\subsection{You only linearize once}

\section{Experiments}

\subsection{Nonlinear Heat Equation}

\begin{equation}
\frac{partial T_i }{\partial t} = \nabla \kappa \nabla T_i = (\kappa_{i+1/2})
\end{equation}

%\subsection{Extreme Learning MLP PINN}

\subsection{CUTEst problems}

\subsection{Pole-cart collocation}

  \begingroup
  \setlength{\arraycolsep}{4pt}
  \renewcommand{\arraystretch}{1.25}
  \begin{table}
  \centering
  \footnotesize
  \begin{tabular}{@{}l c c c c c@{}}
  \toprule
  Operation & Dense & values & in\_cols & start & end \\
  \midrule
  \textit{(seed)}\ $x = \texttt{Identity}(5)$ &
    $\begin{pmatrix*}[r]
      1 & 0 & 0 & 0 & 0 \\ 
      0 & 1 & 0 & 0 & 0 \\
      0 & 0 & 1 & 0 & 0 \\
      0 & 0 & 0 & 1 & 0 \\
      0 & 0 & 0 & 0 & 1
    \end{pmatrix*}$ &
    $\begin{bmatrix}1\\1\\1\\1\\1\end{bmatrix}$ &
    $\begin{bmatrix}0\\1\\2\\3\\4\end{bmatrix}$ &
    $0$ & $5$ \\\addlinespace[8pt]

  $a = c \odot x$,\ $c=(3,5,7,9,11)$ &
    $\begin{pmatrix*}[r]
      3 & 0 & 0 & 0 &  0 \\
      0 & 5 & 0 & 0 &  0 \\
      0 & 0 & 7 & 0 &  0 \\
      0 & 0 & 0 & 9 &  0 \\
      0 & 0 & 0 & 0 & 11
    \end{pmatrix*}$ &
    $\begin{bmatrix}3\\5\\7\\9\\11\end{bmatrix}$ &
    $\begin{bmatrix}0\\1\\2\\3\\4\end{bmatrix}$ &
    $0$ & $5$ \\ \addlinespace[8pt]

  $b = a_{1:5}$ &
    $\begin{pmatrix*}[r]
      0 & 5 & 0 & 0 &  0 \\
      0 & 0 & 7 & 0 &  0 \\
      0 & 0 & 0 & 9 &  0 \\
      0 & 0 & 0 & 0 & 11
    \end{pmatrix*}$ &
    $\begin{bmatrix}5\\7\\9\\11\end{bmatrix}$ &
    $\begin{bmatrix}1\\2\\3\\4\end{bmatrix}$ &
    $0$ & $4$ \\ \addlinespace[8pt]

  $d = a_{0:4}$ &
    $\begin{pmatrix*}[r]
      3 & 0 & 0 & 0 & 0 \\
      0 & 5 & 0 & 0 & 0 \\
      0 & 0 & 7 & 0 & 0 \\
      0 & 0 & 0 & 9 & 0
    \end{pmatrix*}$ &
    $\begin{bmatrix}3\\5\\7\\9\end{bmatrix}$ &
    $\begin{bmatrix}0\\1\\2\\3\end{bmatrix}$ &
    $0$ & $4$ \\ \addlinespace[8pt]

  $e = b - d$ &
    $\begin{pmatrix*}[r]
      -3 & 5 & 0 & 0 &  0 \\
       0 & -5 & 7 & 0 &  0 \\
       0 & 0 & -7 & 9 &  0 \\
       0 & 0 & 0 & -9 & 11
    \end{pmatrix*}$ &
    $\begin{bmatrix}
      -3 & 5 \\
      -5 & 7 \\
      -7 & 9 \\
      -9 & 11
    \end{bmatrix}$ &
    $\begin{bmatrix}
      0 & 1 \\
      1 & 2 \\
      2 & 3 \\
      3 & 4
    \end{bmatrix}$ &
    $0$ & $4$ \\ \addlinespace[8pt]

  $f = \texttt{pad}(e,\,1\!\rightarrow\!0)$ &
    $\begin{pmatrix*}[r]
       0 & 0 & 0 & 0 &  0 \\
      -3 & 5 & 0 & 0 &  0 \\
       0 & -5 & 7 & 0 &  0 \\
       0 & 0 & -7 & 9 &  0 \\
       0 & 0 & 0 & -9 & 11
    \end{pmatrix*}$ &
    $\begin{bmatrix}
      -3 & 5 \\
      -5 & 7 \\
      -7 & 9 \\
      -9 & 11
    \end{bmatrix}$ &
    $\begin{bmatrix}
      0 & 1 \\
      1 & 2 \\
      2 & 3 \\
      3 & 4
    \end{bmatrix}$ &
    $1$ & $5$ \\
  \bottomrule
  \end{tabular}
  \caption{Walk for $f(x) = \mathrm{pad}\bigl((c\odot x)_{1:} - (c\odot x)_{:-1}\bigr)$ with $c=(3,5,7,9,11)$, $\mathrm{len}(x)=5$.
  The mathematical Jacobian (parentheses, right-aligned for sign clarity) is contrasted with the BEllpack
  representation (square brackets, the literal in-memory tensors). Elementwise scaling rewrites
  \texttt{values} in place; slicing shifts the columnfband; subtraction concatenates two $k{=}1$ bands
  into a $k{=}2$ band; pad updates only \texttt{start}/\texttt{end}.}
  \end{table}
  \endgroup

\begin{equation*}
\text{Identity} \rightarrow \text{Constant Diagonal} \rightarrow \text{Diagonal} \rightarrow \text{BEllpack} \rightarrow \text{BCOO}
\end{equation*}

\begin{table}
\centering
\caption{Categorised list of linear JAX primitives supported by Ellpax (trailing \texttt{\_p} dropped).}
\renewcommand{\arraystretch}{1.25}
\begin{tabularx}{0.85\textwidth}{lX}
\toprule
Group & Primitive \\
\midrule
Shape-preserving unary & \texttt{conj, copy, imag, neg, real, cumsum, convert\_element\_type\footnote{cumsum is the only primitive that always densifies}} \\
Shape-affecting unary & \texttt{pad, reshape, broadcast\_in\_dim, reduce\_sum, squeeze}  \\
Indexing & \texttt{slice, rev, gather, transpose, scatter\_add} \\
Shape-preserving multilinear & \texttt{add\footnote{Including add\_any and add\_jaxvals}, sub, mul, div} \\%($+$, $-$, $\times$, $\/$,  $+=$)} \\
Shape-affecting multilinear & \texttt{dot\_general, concatenate} \\ %($@$, $\parallel$) \\
Control flow & \texttt{cond, select\_n, (jit)}\\
Multi-output unary & \texttt{split} \\
\bottomrule
\end{tabularx}
\end{table}

\section{Coloring-pathologic projective planes}

Many of the CUTEst problems involve objective functions of the form:

\begin{equation}
f(x) = \sum_i f_i(x),
\end{equation}

where 




  % \begin{figure}
  % \centering
  % \begin{tikzpicture}[
  %   scale=0.9,
  %   pt/.style={circle, fill=black, inner sep=1.5pt},
  %   lbl/.style={font=\scriptsize},
  % ]
  % % ---- (a) g1 lines (j = s*i + t mod 3) ----
  % \begin{scope}
  %   \node[lbl,above] at (1, 2.55) {(a) $g_1$: $j \equiv s\,i + t$};
  %   \foreach \i in {0,1,2}\foreach \j in {0,1,2} \node[pt] at (\i,\j) {};
  %   % representative lines, one per slope, t=0
  %   \draw[blue!70, thick]   (-0.2,0)   -- (2.2,0);                % s=0
  %   \draw[red!70,  thick]   (-0.2,-0.2)-- (2.2,2.2);              % s=1
  %   \draw[orange!85!black, thick] (0,0) -- (1,2) -- (2,1);        % s=2 (wraps)
  %   \node[lbl] at (1,-0.7) {chrom.\ no.\ $= 3$};
  % \end{scope}

  % % ---- (b) g2 lines (i = s*j + t mod 3) ----
  % \begin{scope}[xshift=4.4cm]
  %   \node[lbl,above] at (1, 2.55) {(b) $g_2$: $i \equiv s\,j + t$};
  %   \foreach \i in {0,1,2}\foreach \j in {0,1,2} \node[pt] at (\i,\j) {};
  %   \draw[blue!70, thick]   (0,-0.2)   -- (0,2.2);
  %   \draw[red!70,  thick]   (-0.2,-0.2)-- (2.2,2.2);
  %   \draw[orange!85!black, thick] (0,0) -- (2,1) -- (1,2);
  %   \node[lbl] at (1,-0.7) {chrom.\ no.\ $= 3$};
  % \end{scope}

  % % ---- (c) union: every pair of points joined by some line ----
  % \begin{scope}[xshift=8.8cm]
  %   \node[lbl,above] at (1, 2.55) {(c) $g_1 + g_2$: $K_9$};
  %   % all 36 pairwise edges, faint
  %   \foreach \i in {0,...,8} {\foreach \k in {0,...,8} {
  %     \ifnum\i<\k
  %       \pgfmathtruncatemacro{\xa}{int(\i/3)} \pgfmathtruncatemacro{\ya}{mod(\i,3)}
  %       \pgfmathtruncatemacro{\xb}{int(\k/3)} \pgfmathtruncatemacro{\yb}{mod(\k,3)}
  %       \draw[gray!35,thin] (\xa,\ya) -- (\xb,\yb);
  %     \fi
  %   }}
  %   % highlight: pair (0,0)-(1,0) covered by g1 only (horizontal)
  %   \draw[red!70,  ultra thick] (0,0) -- (1,0);
  %   % pair (0,0)-(0,1) covered by g2 only (vertical)
  %   \draw[blue!70, ultra thick] (0,0) -- (0,1);
  %   \foreach \i in {0,1,2}\foreach \j in {0,1,2} \node[pt] at (\i,\j) {};
  %   \node[lbl] at (1,-0.7) {chrom.\ no.\ $= 9$};
  % \end{scope}

  % \end{tikzpicture}
  % \caption{Each summand admits a valid 3-coloring of its column-conflict graph
  % (columns sharing $i$ never co-occur in a $g_1$-line; similarly $j$ for $g_2$).
  % Their union makes \emph{every} pair of grid points share a line --- the column
  % conflict graph becomes $K_9$, requiring 9 colors. The highlighted edges in
  % (c) show two pairs each covered by exactly one summand.}
  % \end{figure}

  \begin{figure}
  \centering
  \definecolor{colA}{RGB}{255,182,193}
  \definecolor{colB}{RGB}{173,216,230}
  \definecolor{colC}{RGB}{200,230,180}
  \begin{tikzpicture}[
    scale=0.36,
    lbl/.style={font=\scriptsize},
    sym/.style={font=\Large},
  ]

  % Nonzero patterns at P=3.  Row k = s*P+t, col c = i*P+j.
  \def\gOneRows{0/{0,3,6}, 1/{1,4,7}, 2/{2,5,8},
                3/{0,4,8}, 4/{1,5,6}, 5/{2,3,7},
                6/{0,5,7}, 7/{1,3,8}, 8/{2,4,6}}
  \def\gTwoRows{0/{0,1,2}, 1/{3,4,5}, 2/{6,7,8},
                3/{0,4,8}, 4/{2,3,7}, 5/{1,5,6},
                6/{0,5,7}, 7/{1,3,8}, 8/{2,4,6}}
  \def\gSumRows{0/{0,1,2,3,6}, 1/{1,3,4,5,7}, 2/{2,5,6,7,8},
                3/{0,4,8},     4/{1,2,3,5,6,7}, 5/{1,2,3,5,6,7},
                6/{0,5,7},     7/{1,3,8},     8/{2,4,6}}

  \def\classcolour#1{\ifcase#1 colA\or colB\or colC\fi}

  % --- Panel A: J(g1), columns coloured by  i = c // 3 ---
  \begin{scope}[shift={(0,0)}]
    \foreach \r/\cols in \gOneRows
      \foreach \c in \cols {
        \pgfmathtruncatemacro{\cls}{int(\c/3)}
        \fill[\classcolour{\cls}] (\c,-\r) rectangle ++(1,-1);
      }
    \foreach \r in {0,...,8}\foreach \c in {0,...,8}
      \draw[gray!55, very thin] (\c,-\r) rectangle ++(1,-1);
    \node[lbl] at (4.5,-10.3) {(a) $J(g_1)$, 3-colouring by $i$};
  \end{scope}

  \node[sym] at (10.5, -4.5) {$+$};

  % --- Panel B: J(g2), columns coloured by  j = c mod 3 ---
  \begin{scope}[shift={(12,0)}]
    \foreach \r/\cols in \gTwoRows
      \foreach \c in \cols {
        \pgfmathtruncatemacro{\cls}{mod(\c,3)}
        \fill[\classcolour{\cls}] (\c,-\r) rectangle ++(1,-1);
      }
    \foreach \r in {0,...,8}\foreach \c in {0,...,8}
      \draw[gray!55, very thin] (\c,-\r) rectangle ++(1,-1);
    \node[lbl] at (4.5,-10.3) {(b) $J(g_2)$, 3-colouring by $j$};
  \end{scope}

  \node[sym] at (22.5, -4.5) {$=$};

  % --- Panel C: J(g) = J(g1)+J(g2), no valid coloring -> all black ---
  \begin{scope}[shift={(24,0)}]
    \foreach \r/\cols in \gSumRows
      \foreach \c in \cols
        \fill[black] (\c,-\r) rectangle ++(1,-1);
    \foreach \r in {0,...,8}\foreach \c in {0,...,8}
      \draw[gray!55, very thin] (\c,-\r) rectangle ++(1,-1);
    \node[lbl] at (4.5,-10.3) {(c) $J(g_1{+}g_2)$, no $k{<}9$ colouring};
  \end{scope}

  \end{tikzpicture}
  \caption{Boolean Jacobian patterns at $P=3$. Each summand is $9{\times}9$ with $3$ nonzeros per
  row; the column-conflict graph of $J(g_1)$ admits a $3$-colouring grouped by row index $i$
  (panel a, columns of grid points sharing $i$ never co-occur in a $g_1$-line), and likewise
  $J(g_2)$ admits a $3$-colouring by $j$ (panel b, the transverse partition). When the two
  Jacobians are added, every pair of distinct columns is joined by some $g_1$- or $g_2$-line,
  so the aggregate column-conflict graph is complete and no compression below $9$ colours
  exists (panel c, shown without colouring).}
  \end{figure}


\section{Bib test}
\label{sec:bibtest}

This section exercises the bibliography wiring; it should be removed before submission.

Single citation: foundational sparse Jacobian estimation~\cite{curtispowellreid1974sparse}.
Two-cite: graph-theoretic formulation~\cite{colemanmore1983jacobian, colemanmore1984hessian}.
Three consecutive (compressed): coloring lineage~\cite{curtispowellreid1974sparse,
powelltoint1979sparse, colemanmore1983jacobian}.
Survey reference~\cite{gebremedhin2005color}; recent comparator
work~\cite{hilldalle2025sparser, dallehill2026common}.

Partial-separability~\cite{griewanktoint1982partitioned, conngouldtoint1992lancelot};
multifrontal-assembly analogy~\cite{duffreid1983multifrontal, liu1992multifrontal}.
Experiments on the CUTEst sweep~\cite{gouldorbantoint2014cutest, bongartz1995cute},
bit-exact against the Pearlmutter HVP oracle~\cite{pearlmutter1994fast}.
Background AD and sparse storage~\cite{griewankwalther2008evaluating, baydin2018automatic,
saad2003iterative}; host system~\cite{frostigjohnsonleary2018compiling, jax2018github}.

\section{Conclusions}


\begin{ack}
Use unnumbered first level headings for the acknowledgments. All acknowledgments
go at the end of the paper before the list of references. Moreover, you are required to declare
funding (financial activities supporting the submitted work) and competing interests (related financial activities outside the submitted work).
More information about this disclosure can be found at: \url{https://neurips.cc/Conferences/2026/PaperInformation/FundingDisclosure}.


Do {\bf not} include this section in the anonymized submission, only in the final paper. You can use the \texttt{ack} environment provided in the style file to automatically hide this section in the anonymized submission.
\end{ack}

\bibliographystyle{plainnat}
\bibliography{references}


%%%%%%%%%%%%%%%%%%%%%%%%%%%%%%%%%%%%%%%%%%%%%%%%%%%%%%%%%%%%

\appendix

\section{Technical appendices and supplementary material}
Technical appendices with additional results, figures, graphs, and proofs may be submitted with the paper submission before the full submission deadline (see above). You can upload a ZIP file for videos or code, but do not upload a separate PDF file for the appendix. There is no page limit for the technical appendices.

Note: Think of the appendix as ``optional reading'' for reviewers. The paper must be able to stand alone without the appendix; for example, adding critical experiments that support the main claims to an appendix is inappropriate.

\section{Action of primitive rules on a BEllpack operand}
\label{app:bellpack-rules}

We use the 5-tuple shorthand $\mathrm{BE}[V,\,C,\,s,\,e,\,T]$ for a BEllpack
with values $V$, in-column tuple $C$ of length $k$, row range $[s,e)$, and
transpose flag $T$.  The remaining fields ($\mathtt{out\_size}$,
$\mathtt{in\_size}$, $\mathtt{batch\_shape}$) pass through unchanged unless
explicitly modified.  Each row below is the canonical structural case
implemented in the current source; rules with several distinct branches
are listed once per branch with a labelled left-hand side.  When the
listed branch does not apply, the rule densifies via $V := \mathrm{BE.todense}$
or promotes to BCOO; those fall-back paths are not enumerated here.

\begingroup
\renewcommand{\arraystretch}{1.25}
\footnotesize
\begin{longtable}{@{}p{0.40\textwidth} l@{}}
\toprule
Operation & Result \\
\midrule
\endfirsthead
\toprule
Operation & Result \\
\midrule
\endhead
\bottomrule
\endfoot

\multicolumn{2}{@{}l}{\itshape Shape-preserving unaries} \\
\addlinespace[2pt]
$\Phi(\mathrm{BE})$\footnote{$\Phi$ stands for any element-wise unary that preserves shape and the zero element: $\Phi\in\{\texttt{neg},\texttt{copy},\texttt{conj},\texttt{real},\texttt{imag},\texttt{convert\_element\_type}\}$. Each maps $V\mapsto\Phi(V)$ pointwise on the values tensor, leaving $C$, $s$, $e$, $T$ untouched.} & $\mathrm{BE}[\Phi(V),\,C,\,s,\,e,\,T]$ \\
$\texttt{cumsum}(\mathrm{BE})$ & $\texttt{cumsum}(\mathrm{BE.todense})$ \quad (no structural path) \\
\addlinespace[6pt]

\multicolumn{2}{@{}l}{\itshape Shape-affecting unaries} \\
\addlinespace[2pt]
$\texttt{pad}(\mathrm{BE},\,\mathrm{before},\,\mathrm{after})$, no interior & $\mathrm{BE}[V,\,C,\,s+\mathrm{before},\,e+\mathrm{before},\,T]$ \\
$\texttt{slice}(\mathrm{BE},\,[\ell,h))$, unit stride on out-axis & $\mathrm{BE}[V_{[\ell:h]},\,C_{[\ell:h]},\,s-\ell,\,e-\ell,\,T]$ \\
$\texttt{squeeze}(\mathrm{BE})$, $\mathtt{out\_size}{=}1$ at $s{=}0,e{=}1$ & $\mathrm{BE}[V,\,C,\,s,\,e,\,T]$ \quad (identity) \\
$\texttt{reshape}(\mathrm{BE})$, batched$\to$unbatched, $T{=}\mathrm{False}$ & $\mathrm{BE}[V_{\mathrm{flat}},\,C_{\mathrm{tile}},\,0,\,\prod b\cdot\mathtt{nrows},\,T]$ \\
$\texttt{reshape}(\mathrm{BE})$, unbatched, leading singleton & $\mathrm{BE}[V_{\mathrm{reshape}},\,C_{\mathrm{tile}},\,s,\,e,\,T]$, $b'{=}(1,\dots,1)$ \\
$\texttt{broadcast\_in\_dim}(\mathrm{BE})$, $1$-row $T{=}\mathrm{True}$ tile to $N$ & $\mathrm{BE}[V,\,\mathrm{broadcast}(C,N),\,0,\,N,\,\mathrm{True}]$ \\
$\texttt{broadcast\_in\_dim}(\mathrm{BE})$, leading-batch insert & $\mathrm{BE}[V_{\mathrm{bcast}},\,C_{\mathrm{bcast}},\,s,\,e,\,T]$, $b'{=}(\mathrm{new\_batch}, b)$ \\
$\texttt{reduce\_sum}(\mathrm{BE})$, axis$=$out, static $C$ & $\mathrm{BE}[V_\Sigma,\,\mathrm{dedup}(C),\,0,\,1,\,T]$ \\
$\texttt{reduce\_sum}(\mathrm{BE})$, axis$\subseteq$batch & $\mathrm{BE}[V.\mathrm{sum}(\mathrm{axes}),\,C,\,s,\,e,\,T]$, $b'{=}b\setminus\mathrm{axes}$ \\
\addlinespace[6pt]

\multicolumn{2}{@{}l}{\itshape Indexing} \\
\addlinespace[2pt]
$\texttt{transpose}(\mathrm{BE},\,\mathrm{perm})$, $2{\times}2$ cross-$V$ swap & $\mathrm{BE}[V,\,C,\,s,\,e,\,\neg T]$ \\
$\texttt{transpose}(\mathrm{BE},\,\mathrm{perm})$, batch-only / $V$-fixed & $\mathrm{BE}[V_{\mathrm{perm}},\,C_{\mathrm{perm}},\,s,\,e,\,T]$ \\
$\texttt{rev}(\mathrm{BE})$, axis$=$out, unbatched & $\mathrm{BE}[\mathtt{flip}(V),\,\mathtt{flip}(C),\,N{-}e,\,N{-}s,\,T]$ \\
$\texttt{gather}(\mathrm{BE},\,\dots)$ & raises \texttt{NotImplementedError} (densify upstream) \\
$\texttt{scatter\_add}(\cdot,\,\mathrm{BE},\,\mathrm{idx})$, kept form & $\texttt{scatter\_add\_op}(V,\,\mathrm{idx},\,\dots)$ via dispatch \\
\addlinespace[6pt]

\multicolumn{2}{@{}l}{\itshape Shape-preserving multilinear} \\
\addlinespace[2pt]
$\alpha\cdot\mathrm{BE}$, scalar $\alpha$ & $\mathrm{BE}[\alpha\,V,\,C,\,s,\,e,\,T]$ \\
$\mathbf{w}\odot\mathrm{BE}$, per-row weight broadcast & $\mathrm{BE}[\mathbf{w}\,V,\,\mathrm{broadcast}(C),\,s,\,e,\,T]$ \\
$\mathrm{BE}/\alpha$, scalar $\alpha$ & $\mathrm{BE}[V/\alpha,\,C,\,s,\,e,\,T]$ \\
$\mathrm{BE}_1+\mathrm{BE}_2$, identical $(s,e,b,C)$ & $\mathrm{BE}[V_1+V_2,\,C,\,s,\,e,\,T]$ \\
$\mathrm{BE}_1+\mathrm{BE}_2$, identical $(s,e,b)$, mixed $C$ & $\mathrm{dedup}(\mathrm{BE}[V_1\,\Vert\,V_2,\,C_1\,\Vert\,C_2,\,s,\,e,\,T])$ \\
$\mathrm{BE}_1+\mathrm{BE}_2$, row-disjoint, abutting & $\mathrm{BE}[V_1\,\Vert\,V_2,\,C_1\,\Vert\,C_2,\,\min(s_i),\,\max(e_i),\,T]$ \\
$\mathrm{BE}+\mathrm{BCOO}$ (or any incompatible) & $\mathrm{BCOO}(\mathrm{BE})+\mathrm{BCOO}$ via shared concat path \\
$\mathrm{BE}-\mathrm{BE}'$ & $\mathrm{BE}+(-\mathrm{BE}')$ via the add path \\
\addlinespace[6pt]

\multicolumn{2}{@{}l}{\itshape Shape-affecting multilinear} \\
\addlinespace[2pt]
$\mathrm{BE}\cdot M$, $M$ closure 2-D, contract on in-axis & $\mathrm{dense}\;V_{\mathrm{contract}}$ of shape $(b,N,B)$ where $B{=}M.\mathrm{shape}[1]$ \\
$\texttt{concatenate}(\mathrm{BE}_1,\mathrm{BE}_2)$, axis$<n_{\mathrm{batch}}$ & $\mathrm{BE}[V_1\,\Vert_{\mathrm{ax}}\,V_2,\,C_{\mathrm{widen}},\,s,\,e,\,T]$ \\
$\texttt{concatenate}(\mathrm{BE}_1,\mathrm{BE}_2)$, axis$=$out & $\mathrm{BE}[V_1\,\Vert\,V_2,\,C_1\,\Vert\,C_2,\,0,\,N_1+N_2,\,T]$ \\
\addlinespace[6pt]

\multicolumn{2}{@{}l}{\itshape Control flow / selection} \\
\addlinespace[2pt]
$\texttt{cond}(\mathrm{pred},\dots)$, $\mathrm{pred}$ static & recurse into chosen branch jaxpr \\
$\texttt{jit}(\mathrm{BE})$ & recurse into inner jaxpr; outputs marked traced \\
$\texttt{select\_n}(p,\,\mathrm{BE},\,\mathrm{closures})$, single traced & $\mathrm{BE}[\mathrm{mask}_p\cdot V,\,C,\,s,\,e,\,T]$ \\
$\texttt{select\_n}(p,\,\mathrm{BE}_1,\,\mathrm{BE}_2)$, identical $(s,e,b,C)$ & $\mathrm{BE}[\texttt{select\_n}(p,V_1,V_2),\,C,\,s,\,e,\,T]$ \\
\addlinespace[6pt]

\multicolumn{2}{@{}l}{\itshape Multi-output unary} \\
\addlinespace[2pt]
$\texttt{split}(\mathrm{BE})$, batched, axis$=$out & list of $\mathrm{BE}[V_{[c_i]},\,C_{[c_i]},\,0,\,e_i{-}s_i,\,T]$ \\
$\texttt{split}(\mathrm{BE})$, unbatched, static $C$ & list mixing $\mathrm{BE}$ chunks and empty BCOO chunks \\

\end{longtable}
\endgroup

The rules above cover the BEllpack-preserving paths in
\texttt{lineaxpr/\_linops/ellpack\_transforms.py},
\texttt{lineaxpr/\_linops/ellpack.py}, and
\texttt{lineaxpr/\_rules/\{add,mul,multilinear,structural,control\_flow\}.py}.
Several primitives (\texttt{reshape}, \texttt{broadcast\_in\_dim},
\texttt{reduce\_sum}) carry additional structural branches for less common
shape combinations; we have shown the canonical case for each and refer
the reader to the cited files for the full case-split.

%%%%%%%%%%%%%%%%%%%%%%%%%%%%%%%%%%%%%%%%%%%%%%%%%%%%%%%%%%%%

\newpage
\input{checklist.tex}


\end{document}
